The mineral content in cocoyam leaves is substantially influenced by how processing and storage steps work together, rather than by any single factor. Interaction effects determine both the initial mineral levels and their evolution during storage. Postharvest handling therefore plays a decisive role in micronutrient retention. Storage conditions amplify or reduce nutrient losses, depending on prior processing steps. Among the minerals studied, Phosphorus and calcium seemed to be the most sensitive, showing sharp declines over 42 days, while sodium remained largely unaffected by temperature and duration. 
The observed moisture increase (from about 7.5-7.8\% to 10-13\%) likely contributed to mineral dilution and structural changes in the dried matrix.
Warm storage (32 °C) consistently worsened calcium and phosphorus stability, which may reflect enhanced chemical degradation and mineral redistribution in reactive mineral forms.
The significant interaction between drying method and storage temperature demonstrates the importance of combinating initial drying techniques with appropriate storage protocols to preserve nutrient quality. 
  Overall, the findings suggest that suitable combinations of drying, pretreatment, and controlled storage temperature are essential for maintaining micronutrient quality in dried leafy vegetables. In relation to micronutrient retention in fortified or processed foods, this study highlights that storage duration and temperature can strongly affect mineral stability, particularly when combined with specific processing methods. Practically, the study recommends that processors should focus on low-temperature drying, utilize gentle pretreatments that can reduce mineral loss. Additionally, they should implement short duration storage under cool, dry conditions in order to preserve the nutritional quality \autocite{li5cocoyam}. 
These findings resemble those of \citet{chepkoech}, who observed similar temperature- and humidity-driven losses in vitamins within fortified maize flour. Both studies underline that cool, dry storage conditions are essential to maintain micronutrient integrity.
To compare vitamin and mineral stability, it has been generally reported that vitamin losses are higher than mineral losses in fortified foods during storage \autocite{chepkoech}.
